Working with a variety of digital and physical inputs and outputs (Max/MSP, PlugData, Reaper, TouchDesigner, Houdini, Unity,  ESP32 and Teensy boards, etc.) yet still primarily focusing on the sonic domain, my practice centers on modification (and the often diagrammatic "mappings" that precede the need/desire for modification) as the primary mode of knowledge generation. These three works share common characteristics: all address the metamorphosis and transfiguration of Black being that occurs through space and time, all use sound as their primary expressive medium, all involve material interfaces that demand embodied interaction, and all operate speculatively to refuse dominant historical and technological narratives. Yet each work investigates different material substrates and interaction modalities. \emph{Buffalo Soldier} explores tactility and substrate vibration, while \emph{Did Sun Ra Fear Cryosleep?} investigates fluid dynamics and sedimentation through hydrophone-based sonification, and then \emph{Don't Play with My Hair} examines gestural and structural properties through a custom synthesizer architecture grounded in the materiality of hair.

There is, however, a crucial aspect of this practice that requires explicit acknowledgment: \textbf{both \textit{hardware} modification and \textit{software} modification \textit{are} an embodied practice}. This Cartesian split that positions "software work" as disembodied "mind work," as if coding, patching, and digital manipulation happen in some abstract computational space divorced from physical experience and material understanding, fundamentally misunderstands how knowledge emerges through technological practice. My body is fully engaged when I modify code or build Max patches in much the same way it is when I patch together physical computing components: my eyes track data visualizations and waveform displays, my ears evaluate subtle sonic shifts across iterative adjustments, my hands manipulate materials and route signals. My posture shifts during extended sessions of experimentation, I stand, I pace. I sketch both physically and mentally. My entire sensorium registers the friction between what I want to hear or see and what the tools defaults initially afford. 

Modification emerges from sensory feedback loops---listen, adjust, listen again, adjust again---decisions guided by what Manning and Massumi call "thinking-feeling," the inseparability of conceptual work from bodily, sensorial engagement \parencite[41-42]{Manning_Massumi_2014}. When I route extracted hair data (tables of values) through a modified wavetable object, the decision about which configuration "feels right" arises from my body registering alignment between sonic output and the cultural memory I am attempting to encode. This is embodied knowledge production, and the works presented in this chapter emerge from thousands of such micro-decisions made through sustained physical engagement with both computational and physical materials (from the virtual to the actualized). To dismiss this as mere "computer work" is to reproduce some of the very assumptions that ST seeks to unsettle. To say that only certain forms of material manipulation constitute legitimate embodied practice while computational practice remains abstractly intellectual is itself an epistemological fallacy.